% !TEX root = main.tex
\section{第一章：代数结构}
本章选取作业中的19题。题目及证明思路依据现有笔记的习题文件和所引用的正文整理；下列细节是沿用这些思路的补充推导，不标作教材原证。各题“原文依据”指实际核对的现有笔记，未核对到教材证明的部分不声称已与教材原证核对。

\paragraph{本章常用记号与判别方法}
群的单位元记为$e$。子群$H\leq G$的左、右陪集分别为$aH$、$Ha$；左陪集个数记为指数$[G:H]$。若所有$a\in G$均满足$aH=Ha$，则称$H$为正规子群，记为$H\trianglelefteq G$。群中有消去律，但一般不能交换因子的次序。
非空子集$H$为子群的常用判别是：任意$x,y\in H$均有$xy^{-1}\in H$。环的双边理想$I$是加法子群，并且满足$RI\subseteq I$、$IR\subseteq I$。下文的“理想”均指双边理想。幂零元是某个正整数次幂为零的元素；零因子按笔记约定必须是非零元素。

\subsection{习题1：平方公式判别交换群}
\paragraph{题目}
若群$G$中任意$a,b$满足$(ab)^2=a^2b^2$，证明$G$交换。
\paragraph{考点与思路}
将幂展开，只消去两端的因子。也可证明该条件与交换性等价。
\textbf{原文依据：}现有笔记本题解答；群的消去律。
\paragraph{所用结论}
\toolref{g-laws}。
\begin{proof}[详细证明]
任取$a,b\in G$，题设给出$abab=aabb$。先在两端左乘$a^{-1}$，再右乘$b^{-1}$，得到
\begin{align*}
abab&=aabb,\\
bab&=abb,\\
ba&=ab.
\end{align*}
因为$a,b$任意，群中任意两元素交换，故$G$为交换群。

反过来，若$G$交换，则$abab=aabb=a^2b^2$，所以题设条件也成立。
\end{proof}
\paragraph{复习提示}
消去律仅允许消去位于同一侧的相同因子；不能从一开始就把中间的$ba$换成$ab$。

\subsection{习题2：每个元素的平方为单位元}
\paragraph{题目}
若群$G$中的每个元素$a$都满足$a^2=e$，证明$G$交换。
\paragraph{考点与思路}
将题设同时用于$a$、$b$和$ab$，再使用乘积取逆时的反序公式。
\textbf{原文依据：}现有笔记本题解答；乘积取逆公式。
\paragraph{所用结论}
\toolref{g-laws}。
\begin{proof}[详细证明]
由$a^2=e$知$a$本身就是$a$的逆元，所以每个$a\in G$都有$a^{-1}=a$。
任取$a,b\in G$，因为$ab\in G$，它也满足$(ab)^{-1}=ab$。另一方面，任何群中都有$(ab)^{-1}=b^{-1}a^{-1}$，于是
\[
ab=(ab)^{-1}=b^{-1}a^{-1}=ba.
\]
因此$G$交换。
\end{proof}
\paragraph{复习提示}
题设包含单位元，并不意味着每个元素的阶都恰为2；单位元的阶为1。结论也不要求$G$有限。

\subsection{习题4：可解方程的结合系统}
\paragraph{题目}
非空集合$G$上有一个结合的二元运算。若对任意$a,b\in G$，方程$ax=b$与$ya=b$都有解，证明$G$构成群。
\paragraph{考点与思路}
尚未证明是群时，不能使用逆元和消去律。先将某个元素的单侧单位元推广为全体元素的单侧单位元。
\textbf{原文依据：}现有笔记本题解答中先构造双边单位元、再构造双边逆元的步骤。
\paragraph{所用结论}
\toolref{g-group-construction}。
\begin{proof}[详细证明]
\textbf{第一步：构造双边单位元。}
因为$G$非空，可固定$a\in G$。由题设，存在$e,f\in G$使
\[
ae=a,\qquad fa=a.
\]
任取$x\in G$。方程$ya=x$有解，因此$x=ya$，结合律给出
\[
xe=(ya)e=y(ae)=ya=x.
\]
所以$e$是整个集合的右单位元。又由方程$az=x$有解，可写$x=az$，从而
\[
fx=f(az)=(fa)z=az=x.
\]
所以$f$是整个集合的左单位元。将两者用于$fe$，分别得到$fe=f$与$fe=e$，故$f=e$。因此$e$为双边单位元。

\textbf{第二步：构造双边逆元。}
任取$g\in G$，由方程可解性，存在$b,c\in G$使$gb=e$、$cg=e$。于是
\[
c=ce=c(gb)=(cg)b=eb=b.
\]
所以$gb=bg=e$，即$g$有双边逆元。

二元运算已经具有封闭性，结合律由题设保证；现在单位元和每个元素的逆元也已得到，故$G$是群。
\end{proof}
\paragraph{复习提示}
“有解”不必加强为“唯一解”。在构造出群以前，任何一步都不应预先使用消去律。

\subsection{习题6：任意偶数阶的非交换例子}
\paragraph{题目}
对整数$n>2$，构造一个$2n$阶的非交换群。
\paragraph{考点与思路}
沿用原解答，取正$n$边形的几何对称群，再按正文中相邻顶点像的计数和旋转、反射关系证明其阶数与非交换性。
\textbf{原文依据：}现有笔记本题解答；正文“图形的对称群与二面体群”中对称群命题、二面体群定义及阶数与非交换性的命题。
\paragraph{所用结论}
\toolref{g-dihedral}。
\begin{proof}[详细证明]
取平面上的正$n$边形，记$D_n$为保持此正多边形整体不变的所有平面等距变换的集合。恒等变换属于$D_n$；两个对称变换的复合仍保持图形；逆变换也保持图形。映射复合具有结合律，所以$D_n$在复合下成群。

\textbf{先证明$|D_n|=2n$。}任意对称变换置换顶点，因而保持顶点的重心，即正多边形的中心$O$。固定两个相邻顶点$P,Q$，由于$n>2$，$O,P,Q$不共线。平面等距变换由三个不共线点的像唯一决定，因此一个对称变换由$P,Q$的像唯一决定。
$P$的像有$n$种选择；$Q$的像必须与$P$的像相邻，而每个顶点恰有两个相邻顶点，所以$Q$的像至多有两种选择。故$|D_n|\leq2n$。

令$T$为绕$O$旋转$2\pi/n$，$S$为关于一条固定对称轴的反射。变换
\[
T^k,\quad ST^k\qquad(0\leq k<n)
\]
都保持正多边形。$n$个旋转$T^k$的旋转角在模$2\pi$意义下不同，故彼此不同；左复合$S$可逆，故$n$个$ST^k$也彼此不同。旋转保持平面定向，反射与旋转的复合反转平面定向，所以两类不相交。已经得到$2n$个不同对称变换，结合前面的上界，得$|D_n|=2n$。

\textbf{再证明它不交换。}反射反转旋转方向，所以$STS=T^{-1}$，且$S^2=e$。若$ST=TS$，则右乘$S$得$STS=T$，与前式比较可得$T=T^{-1}$，即$T^2=e$。
但是$T$的阶恰为$n$：旋转$k$次恢复原位置当且仅当$2k\pi/n$为$2\pi$的整数倍，最小正整数$k$为$n$。所以$T^2=e$将给出$n\mid2$，与$n>2$矛盾。故$ST\ne TS$，$D_n$为所需的非交换群。
\end{proof}
\paragraph{复习提示}
本笔记约定$|D_n|=2n$。条件$n>2$不可略去，2阶和4阶群均交换。

\subsection{习题7：反复共轭}
\paragraph{题目}
若群$G$中的$a,b$满足$a^{-1}ba=b^r$，其中$r$为正整数，证明对任意正整数$i$都有$a^{-i}ba^i=b^{r^i}$。
\paragraph{考点与思路}
先验证共轭保持正整数次幂，再对共轭次数$i$归纳。
\textbf{原文依据：}现有笔记本题解答中的共轭幂公式与数学归纳法。
\paragraph{所用结论}
\toolref{g-laws}、\toolref{n-induction}。
\begin{proof}[详细证明]
对任意正整数$m$，将相邻的$aa^{-1}$消去，得
\[
(a^{-1}ba)^m=a^{-1}b^ma.
\]
因此$a^{-1}b^ma=(b^r)^m=b^{rm}$。

当$i=1$时，结论就是题设。假设某个$i\geq1$已有$a^{-i}ba^i=b^{r^i}$，则
\begin{align*}
a^{-(i+1)}ba^{i+1}
&=a^{-1}(a^{-i}ba^i)a\\
&=a^{-1}b^{r^i}a\\
&=b^{r\cdot r^i}=b^{r^{i+1}}.
\end{align*}
归纳法给出所有正整数$i$的结论。
\end{proof}
\paragraph{复习提示}
反复共轭使指数每次乘以$r$，最终指数为$r^i$，不是$ri$。

\subsection{习题8：取逆映射的同态性}
\paragraph{题目}
证明：群$G$交换，当且仅当映射$\iota:G\to G$、$\iota(a)=a^{-1}$是群同态。取逆映射本身为双射，所以本题中的同态条件也等价于它是同构。
\paragraph{考点与思路}
取逆映射总为双射；真正需要判别的是它是否保乘法。
\textbf{原文依据：}现有笔记本题解答中先验证取逆为双射、再比较保乘法与反序公式的步骤。
\paragraph{所用结论}
\toolref{g-laws}、\toolref{g-first-iso}。
\begin{proof}[详细证明]
因为$\iota(\iota(a))=(a^{-1})^{-1}=a$，所以$\iota$的逆映射就是自身，因而总是双射。

\textbf{先由同态推出交换性。}若$\iota$为同态，则保乘法给出
\[
a^{-1}b^{-1}=\iota(a)\iota(b)=\iota(ab)=b^{-1}a^{-1}.
\]
两边取逆得到$ba=ab$。因$a,b$任意，$G$交换。

\textbf{再由交换性推出同构。}若$G$交换，则对任意$a,b\in G$，
\[
\iota(ab)=(ab)^{-1}=b^{-1}a^{-1}=a^{-1}b^{-1}=\iota(a)\iota(b).
\]
所以$\iota$为同态，再结合双射性，知它是同构。
\end{proof}
\paragraph{复习提示}
任何群都满足$(ab)^{-1}=b^{-1}a^{-1}$；保乘法所要求的却是$(ab)^{-1}=a^{-1}b^{-1}$。

\subsection{习题9：由子集定义等价关系}
\paragraph{题目}
设$S\subseteq G$非空，在$G$上定义$a\sim b\iff ab^{-1}\in S$。证明该关系为等价关系，当且仅当$S$为子群。
\paragraph{考点与思路}
等价关系要求自反、对称、传递。反向将关系与单位元联系起来，导出子群判别条件。
\textbf{原文依据：}现有笔记本题解答；子群判别法与右陪集。
\paragraph{所用结论}
\toolref{g-subgroup}、\toolref{g-cosets}。
\begin{proof}[详细证明]
\textbf{若$S$为子群。}任取$a\in G$，因$aa^{-1}=e\in S$，有$a\sim a$，满足自反性。
若$a\sim b$，则$ab^{-1}\in S$，子群对逆封闭给出$ba^{-1}=(ab^{-1})^{-1}\in S$，所以$b\sim a$。
若$a\sim b$且$b\sim c$，则$ab^{-1},bc^{-1}\in S$，因而
\[
ac^{-1}=(ab^{-1})(bc^{-1})\in S,
\]
即$a\sim c$，满足传递性。因此该关系是等价关系。

\textbf{若该关系为等价关系。}自反性给出$e=aa^{-1}\in S$。任取$u,v\in S$，有$u\sim e$、$v\sim e$。由对称性得$e\sim v$，再由传递性得$u\sim v$，故$uv^{-1}\in S$。

说明此条件为何足够：取$u=e$可知$v^{-1}\in S$；再将已经得到的$v^{-1}$代入上述差式条件，得到$u(v^{-1})^{-1}=uv\in S$。所以$S$含单位元并对逆和乘积封闭，确为子群。

此外，$a\sim b$等价于$a=sb$对某个$s\in S$成立，因此$b$所在的等价类为右陪集$Sb$。
\end{proof}
\paragraph{复习提示}
$ab^{-1}$对应右陪集$Sb$；不要与$a^{-1}b$所对应的左陪集混淆。

\subsection{习题10：整数加法群与其倍数子群}
\paragraph{题目}
设$n$为正整数，证明加法群$n\Z$与$\Z$同构。
\paragraph{考点与思路}
同构必须写出映射，并分别验证保运算、单射和满射。
\textbf{原文依据：}现有笔记本题解答中的映射$k\mapsto nk$。
\paragraph{所用结论}
\toolref{g-first-iso}、\toolref{r-quotients}（其中$\Z$为整环，非零整数可消去）。
\begin{proof}[详细证明]
定义$\varphi:\Z\to n\Z$为$\varphi(k)=nk$。对任意$k,l\in\Z$，
\[
\varphi(k+l)=n(k+l)=nk+nl=\varphi(k)+\varphi(l),
\]
所以它保加法。若$\varphi(k)=\varphi(l)$，则$n(k-l)=0$。因为$n\ne0$，整数乘法无零因子，故$k=l$，所以它为单射。
任取$n\Z$的元素，按定义可写成$nk$，它就是$\varphi(k)$，所以映射满射。因此$\Z\cong n\Z$。
\end{proof}
\paragraph{复习提示}
这里讨论加法群。若$n>1$，上述映射不保通常的乘法，例如$\varphi(1\cdot1)=n$，而$\varphi(1)\varphi(1)=n^2$。若允许$n=0$，同构结论不成立。

\subsection{习题12：指数2的子群正规}
\paragraph{题目}
若$G$为$2n$阶群，其中有一个$n$阶子群$H$，证明$H$正规。原题的条件给出$[G:H]=2$；以下按现有笔记的左右陪集方法证明这一更一般的指数2结论。
\paragraph{考点与思路}
指数2意味着除$H$以外，只有一个左陪集和一个右陪集，它们都必须是$H$的补集。
\textbf{原文依据：}现有笔记本题解答中按$a\in H$与$a\notin H$分类、比较左右陪集的证明。
\paragraph{所用结论}
\toolref{g-cosets}、\toolref{g-normal}、\toolref{g-lagrange}。
\begin{proof}[详细证明]
任取$a\in G$。若$a\in H$，子群性质给出$aH=H=Ha$。

若$a\notin H$，则$aH\ne H$，并且不同左陪集不交。因为恰有两个左陪集，必有$aH=G\setminus H$。
取逆映射将左陪集$xH$一一对应到右陪集$Hx^{-1}$，所以右陪集也恰有两个。同理$Ha=G\setminus H$。于是这一情形下也有$aH=Ha$。

对所有$a\in G$均已证明$aH=Ha$，故$H$正规。有限群的特殊情形由$[G:H]=|G|/|H|=2$立即得到。
\end{proof}
\paragraph{复习提示}
一般结论不要求群有限；只有用阶数比计算指数时才使用有限性。

\subsection{习题18：只有有限多个子群的群必有限}
\paragraph{题目}
证明：如果群$G$只有有限多个子群，那么$G$必为有限群。
\paragraph{考点与思路}
先排除无限阶元素，再将全群表示为有限多个有限循环子群的并。
\textbf{原文依据：}现有笔记本题解答中先排除无限阶元素、再取有限多个循环子群之并的证明。
\paragraph{所用结论}
\toolref{g-cyclic}。
\begin{proof}[详细证明]
\textbf{第一步：每个元素都为有限阶。}
反设$a\in G$为无限阶，则子群$\langle a^m\rangle$（$m=1,2,\ldots$）两两不同。
事实上，若$\langle a^m\rangle=\langle a^n\rangle$，则存在整数$u,v$使$a^m=a^{nu}$、$a^n=a^{mv}$。无限阶意味着相等幂的指数相等，故$m=nu$、$n=mv$。由于$m,n$为正整数，$u,v$均为正整数；于是$uv=1$，得$u=v=1$、$m=n$。
这产生无限多个子群，与假设矛盾。

\textbf{第二步：全群有限。}
所有循环子群均为$G$的子群，所以它们只有有限多个，记为$C_1,\ldots,C_t$。第一步说明每个$C_j$都有限。
任意$g\in G$属于其生成的循环子群$\langle g\rangle$，因此
\[
G=\bigcup_{j=1}^{t}C_j,\qquad |G|\leq\sum_{j=1}^{t}|C_j|<\infty.
\]
所以$G$有限。
\end{proof}
\paragraph{复习提示}
仅证明“每个元素有限阶”还不足以证明群有限；必须利用循环子群的个数也有限。

\subsection{习题21：有限子群乘积的计数}
\paragraph{题目}
设$H,K$为群$G$的有限子群，证明
\[
|HK|=\frac{|H||K|}{|H\cap K|}.
\]
此处不要求集合$HK$是子群。
\paragraph{考点与思路}
映射$(h,k)\mapsto hk$一般不单射；准确计算每个像的原像个数。
\textbf{原文依据：}现有笔记本题解答中的满射$H\times K\to HK$及每个像的原像计数。
\paragraph{所用结论}
\toolref{g-laws}、\toolref{g-fiber-count}。
\begin{proof}[详细证明]
定义满射$\pi:H\times K\to HK$，$\pi(h,k)=hk$。固定$g=hk\in HK$的一种表示。

若$h'k'=hk$，两边左乘$h^{-1}$、右乘$k'^{-1}$，得到
\[
t:=h^{-1}h'=kk'^{-1}\in H\cap K.
\]
由此$h'=ht$、$k'=t^{-1}k$。反之，任取$t\in H\cap K$，都有$ht\in H$、$t^{-1}k\in K$，且
\[
(ht)(t^{-1}k)=hk=g.
\]
不同$t$给出不同的有序对，因为$ht_1=ht_2$可消去$h$而得到$t_1=t_2$。因此$g$的全部原像与$H\cap K$一一对应，恰有$|H\cap K|$个。

对有限集合$H\times K$按原像分组计数，得
\[
|H||K|=|HK|\,|H\cap K|,
\]
除以非零正整数$|H\cap K|$即得所求。
\end{proof}
\paragraph{复习提示}
需要扣除的是交集造成的重复表示；即使$HK$不是子群，集合计数公式仍成立。

\subsection{习题22：正规子群的乘积}
\paragraph{题目}
设$M,N\trianglelefteq G$，证明$MN=NM$且$MN\trianglelefteq G$。若$M\cap N=\{e\}$，证明$MN/N\cong M$。
\paragraph{考点与思路}
先证明集合乘积是正规子群，再构造商群的满同态并求核。
\textbf{原文依据：}现有笔记本题解答中先证乘积相等与正规、再构造$m\mapsto mN$的步骤；同态基本定理。
\paragraph{所用结论}
\toolref{g-subgroup}、\toolref{g-normal}、\toolref{g-first-iso}。
\begin{proof}[详细证明]
\textbf{第一步：证明$MN=NM$。}
任取$m\in M,n\in N$，因$M$正规，$n^{-1}mn\in M$，于是$mn=n(n^{-1}mn)\in NM$，故$MN\subseteq NM$。
同样，由$M$正规，$nmn^{-1}\in M$，所以$nm=(nmn^{-1})n\in MN$，得到反向包含。

\textbf{第二步：证明$MN$为正规子群。}
$e\in MN$。集合乘积满足结合律，利用$NM=MN$，得
\[
(MN)(MN)=M(NM)N=M(MN)N=MN,
\qquad (MN)^{-1}=NM=MN.
\]
所以$MN$对乘法和逆封闭，为子群。再任取$g\in G$，由两个子群的正规性，
\[
g(MN)g^{-1}=(gMg^{-1})(gNg^{-1})=MN.
\]
因此$MN\trianglelefteq G$。

\textbf{第三步：构造同构。}
因为$N$在$G$中正规，它在子群$MN$中也正规，可以构成商群$MN/N$。定义
\[
\varphi:M\longrightarrow MN/N,\qquad \varphi(m)=mN.
\]
商群运算给出$\varphi(m_1m_2)=m_1m_2N=(m_1N)(m_2N)$，所以$\varphi$是同态。
任意陪集$mnN$等于$mN$，故$\varphi$满射。其核为
\[
\ker\varphi=\{m\in M:mN=N\}=M\cap N.
\]
因此同态基本定理给出$M/(M\cap N)\cong MN/N$。当$M\cap N=\{e\}$时，$\varphi$本身单射，故$M\cong MN/N$。
\end{proof}
\paragraph{复习提示}
证明商群同构时，应逐项交代映射、保运算、满射和核；交平凡正是单射的依据。

\subsection{习题29：改变坐标得到同构环}
\paragraph{题目}
在交换幺环$R$的底集上定义
\[
a\oplus b=a+b-1,\qquad a\odot b=a+b-ab.
\]
证明新运算使它成为与原环同构的交换幺环，并求新的零元、单位元及加法逆元。
\paragraph{考点与思路}
寻找坐标变换$f(a)=1-a$，把新运算转回熟悉的运算。
\textbf{原文依据：}现有笔记本题解答中的坐标变换$f(a)=1-a$。
\paragraph{所用结论}
\toolref{g-transport}。
\begin{proof}[详细证明]
记具有新运算的集合为$R_{\mathrm{new}}$，定义$f:R_{\mathrm{new}}\to R$为$f(a)=1-a$。因为$f(f(a))=1-(1-a)=a$，它是双射。
对任意$a,b\in R$，
\begin{align*}
f(a\oplus b)&=1-(a+b-1)=2-a-b=f(a)+f(b),\\
f(a\odot b)&=1-(a+b-ab)=1-a-b+ab\\
&=(1-a)(1-b)=f(a)f(b).
\end{align*}
所以$f$将两种新运算分别对应到原环的加法、乘法。由于$f$单射，原环的每一条运算恒等式均可通过$f$拉回；例如
\begin{align*}
f((a\odot b)\odot c)&=(f(a)f(b))f(c)\\
&=f(a)(f(b)f(c))=f(a\odot(b\odot c)),
\end{align*}
从而新乘法满足结合律。加法的结合、交换以及分配律同理由原环的相应公理得到。因此新结构为环，且$f$为环同构。

新零元是$f^{-1}(0)=1$，新乘法单位元是$f^{-1}(1)=0$。也可直接核对$a\oplus1=a$、$a\odot0=a$。元素$a$的新加法逆元为$2-a$，因为
\[
a\oplus(2-a)=a+(2-a)-1=1,
\]
右端正是新零元。新乘法交换性由$a+b-ab=b+a-ba$得到。
\end{proof}
\paragraph{复习提示}
式中$0,1,2$表示原环元素，$2=1+1$；新结构的零元与单位元分别是原来的$1$与$0$。

\subsection{习题30：子环与母环的单位元和交换性}
\paragraph{题目}
分别举例说明：(i)母环有单位元而子环无单位元；(ii)母环无单位元而子环有单位元；(iii)母环和子环都有单位元，但二者不同；(iv)母环不交换而子环交换。
\paragraph{考点与思路}
这里子环只要求对加法、取负和乘法封闭，并不要求与母环共用单位元。应检查实例确实满足对应性质。
\textbf{原文依据：}现有笔记本题解答的四个反例构造及原有顺序。
\paragraph{所用结论}
\toolref{g-ring-criteria}。
\begin{proof}[详细证明]
\textbf{(i)}取$R=\Z$、$S=2\Z$。偶数对差和积封闭，所以$S$是子环；$R$的单位元为1。若$u\in2\Z$是$S$的单位元，则$u\cdot2=2$，在整数中得$u=1$，与$u$为偶数矛盾。因此$S$没有单位元。

\textbf{(ii)}取逐分量运算的环$R=2\Z\times\Z$，以及子环$S=\{0\}\times\Z$。若$R$有单位元$(u,v)$，将它乘以$(2,0)$得到$(2u,0)=(2,0)$，所以$u=1\notin2\Z$，矛盾。而对任意$(0,k)\in S$，
\[
(0,1)(0,k)=(0,k)=(0,k)(0,1),
\]
所以$S$有单位元$(0,1)$。

\textbf{(iii)}取$R=M_2(\R)$，令
\[
S=\left\{\begin{pmatrix}a&0\\0&0\end{pmatrix}:a\in\R\right\}.
\]
两个这样的矩阵之差和积仍有相同形状，所以$S$是子环。母环的单位元为$E=\operatorname{diag}(1,1)$，子环的单位元为$\operatorname{diag}(1,0)$，两者不同。

\textbf{(iv)}仍取$R=M_2(\R)$，令$S=\{aE:a\in\R\}$为数量矩阵构成的子环。因$(aE)(bE)=abE=baE$，$S$交换。记$E_{ij}$为第$(i,j)$项为1、其余项为0的矩阵，则
\[
E_{12}E_{21}=E_{11},\qquad E_{21}E_{12}=E_{22},
\]
二者不同，故母环不交换。
\end{proof}
\paragraph{复习提示}
单位元是相对于环内全部元素定义的；子环中的单位元不必对母环其他元素起单位作用。

\subsection{习题32：域没有非平凡双边理想}
\paragraph{题目}
证明域$F$的双边理想只有$\{0\}$和$F$。同一结论也适用于体。
\paragraph{考点与思路}
非零理想中取非零元素，用其逆元推出单位元属于该理想。
\textbf{原文依据：}现有笔记本题解答中由非零理想元素推出$1\in I$的证明。
\paragraph{所用结论}
\toolref{g-ring-criteria}。
\begin{proof}[详细证明]
设$I$为$F$的理想。若$I=\{0\}$，已经属于所列情形。若$I\ne\{0\}$，可取$a\in I$、$a\ne0$。域中$a$有逆元$a^{-1}\in F$，理想的吸收性给出$1=a^{-1}a\in I$。
任取$r\in F$，再次由吸收性，$r=r\cdot1\in I$，所以$F\subseteq I$。结合$I\subseteq F$，得到$I=F$。

证明只使用非零元素有双边逆元，没有使用乘法交换性，所以对不一定交换的体也成立。
\end{proof}
\paragraph{复习提示}
一般环中非零元素未必可逆，不能对任意环照搬这一步。

\subsection{习题34：理想的根仍为理想}
\paragraph{题目}
设$R$为交换环、$I$为其理想，定义
\[
\sqrt I=\{r\in R:\text{存在正整数 }n\text{ 使 }r^n\in I\}.
\]
证明$\sqrt I$为$R$的理想。
\paragraph{考点与思路}
验证差封闭和吸收性。两个元素对应的幂次数可能不同，统一选取$n+m-1$次幂。
\textbf{原文依据：}现有笔记本题解答中的$(r-s)^{n+m-1}$二项式展开和吸收性验证。
\paragraph{所用结论}
\toolref{g-ring-criteria}、\toolref{g-binomial}。
\begin{proof}[详细证明]
首先$0^1=0\in I$，故$0\in\sqrt I$，集合非空。

任取$r,s\in\sqrt I$，由定义可选正整数$n,m$使$r^n\in I$、$s^m\in I$。令$N=n+m-1$。因为环交换，可以使用二项式展开：
\[
(r-s)^N=\sum_{k=0}^{N}(-1)^{N-k}\binom Nk r^k s^{N-k}.
\]
对任一项，若$k\geq n$，它含有因子$r^n$，故属于$I$；若$k<n$，则$k\leq n-1$，从而$N-k\geq m$，这一项含有因子$s^m$，也属于$I$。整数倍与有限和仍在加法子群$I$中，故$(r-s)^N\in I$，即$r-s\in\sqrt I$。因此$\sqrt I$为加法子群。

任取$t\in R$、$r\in\sqrt I$，选$n\geq1$使$r^n\in I$。由交换性和理想吸收性，
\[
(tr)^n=t^nr^n\in I.
\]
所以$tr\in\sqrt I$；交换性又给出$rt=tr\in\sqrt I$。因此$\sqrt I$满足双边吸收性，确为理想。

若$R$没有单位元，上述展开仍然成立：在$k=0$或$k=N$的端项中，直接写作$(-s)^N$或$r^N$，不引入$R$内不存在的零次幂。判断因子属于$I$时，若剩余因子个数为零，就直接使用$r^n\in I$或$s^m\in I$。
\end{proof}
\paragraph{复习提示}
不能仅由$r^n,s^m\in I$断言$(r-s)^{\max(n,m)}\in I$；$n+m-1$保证展开中每一项都含足够多的某一种因子。

\subsection{习题39：交换环的幂零元构成理想}
\paragraph{题目}
证明交换环中全体幂零元构成理想，并说明非交换环中该结论不成立。
\paragraph{考点与思路}
它是习题34中$I=\{0\}$的特例；复习时也应能独立写出幂次数的论证。
\textbf{原文依据：}现有笔记本题解答；习题34取$I=\{0\}$的情形，以及矩阵单位$E_{12},E_{21}$的非交换反例。
\paragraph{所用结论}
\toolref{g-ring-criteria}、\toolref{g-binomial}。
\begin{proof}[详细证明]
记$\mathcal N=\{a\in R:\text{存在正整数 }n\text{ 使 }a^n=0\}$。因为$0\in\mathcal N$，它非空。
若$a,b\in\mathcal N$，选$n,m\geq1$使$a^n=b^m=0$。展开$(a-b)^{n+m-1}$，其中任意一项$a^k b^{n+m-1-k}$都满足$k\geq n$或$n+m-1-k\geq m$，因而为零。所以
\[
(a-b)^{n+m-1}=0,
\]
即$a-b\in\mathcal N$。若$r\in R$，交换性给出$(ra)^n=r^na^n=0$，所以$ra=ar\in\mathcal N$。因此$\mathcal N$为理想。

非交换环的反例取$M_2(\Q)$中的
\[
A=\begin{pmatrix}0&1\\0&0\end{pmatrix},\qquad
B=\begin{pmatrix}0&0\\1&0\end{pmatrix}.
\]
有$A^2=B^2=0$，但$(A+B)^2=E$。若$A+B$幂零，某个正整数$k$应满足$(A+B)^k=0$，于是$(A+B)^{2k}=0$；另一方面$(A+B)^{2k}=E$，矛盾。因此幂零元之和未必幂零，幂零元集甚至不一定对加法封闭。
\end{proof}
\paragraph{复习提示}
证明中交换性的用途是二项式展开以及$(ra)^n=r^na^n$，不能把这两步用于任意非交换环。

\subsection{习题43：连续函数环中的理想与零因子}
\paragraph{题目}
记$C[0,1]$为$[0,1]$上全部实值连续函数构成的环，加法和乘法逐点定义。
\begin{enumerate}
\item 每个真理想$I\ne C[0,1]$都存在公共零点，即存在$t_0\in[0,1]$使所有$f\in I$均满足$f(t_0)=0$。
\item 对非零函数$f\in C[0,1]$，证明$f$为零因子，当且仅当其零点集包含一个非空的相对开区间。
\end{enumerate}
相对开区间指$(\alpha,\beta)\cap[0,1]$形式的非空集合，其中$\alpha<\beta$为实数。这里“非零函数”指不恒为零的函数。
\paragraph{考点与思路}
第一问用紧致性选有限子覆盖，再以平方和构造处处非零的理想元素；第二问使用连续函数在非零点附近仍非零的性质。
\textbf{原文依据：}现有笔记本题解答中紧致性、有限平方和的构造，以及非零邻域和三角形函数的证明。
\paragraph{所用结论}
\toolref{g-ring-criteria}、\toolref{g-analysis}。
\begin{proof}[详细证明]
\textbf{(i) 真理想具有公共零点。}
反设$I$没有公共零点，则对每个$t\in[0,1]$，都可选$f_t\in I$使$f_t(t)\ne0$。
由连续性存在$t$的相对开邻域$U_t$，使$f_t$在$U_t$上处处非零。这些邻域覆盖$[0,1]$。

闭区间$[0,1]$紧致，故可选有限子覆盖$U_{t_1},\ldots,U_{t_m}$。记$f_j=f_{t_j}$，并令
\[
g=f_1^2+\cdots+f_m^2.
\]
因为$f_j\in I$，理想吸收性给出$f_j^2\in I$，加法封闭性给出$g\in I$。
任意$x\in[0,1]$至少属于某个$U_{t_j}$，因此对应$f_j(x)\ne0$，从而$g(x)=\sum_{j=1}^m f_j(x)^2>0$。
所以函数$h(x)=1/g(x)$在整个区间连续，属于$C[0,1]$。

由理想吸收性，常值函数$1=gh\in I$。于是任取$u\in C[0,1]$都有$u=u\cdot1\in I$，这说明$I=C[0,1]$，与$I$为真理想矛盾。故必有公共零点。

\textbf{(ii) 零因子的必要性。}
设$f\ne0$为零因子，则存在不恒为零的$g\in C[0,1]$使$fg=0$。选$t\in[0,1]$满足$g(t)\ne0$。
连续性给出某个$\delta>0$，使得当$x\in[0,1]$且$|x-t|<\delta$时，
\[
|g(x)-g(t)|<\frac{|g(t)|}{2},\qquad
|g(x)|>\frac{|g(t)|}{2}>0.
\]
在非空相对开区间$(t-\delta,t+\delta)\cap[0,1]$上，$f(x)g(x)=0$且$g(x)\ne0$，所以$f(x)=0$。因此$f$的零点集包含这样的区间。

\textbf{(ii) 零因子的充分性。}
反过来，设$f$在非空相对开区间$J$上恒为零。这样的$J$含有一个位于$(0,1)$内的小闭区间$[c-\varepsilon,c+\varepsilon]$，其中$\varepsilon>0$。这是因为非空相对开集即使包含端点，也包含相邻的内部点，故可在其内部选$c$及足够小的$\varepsilon$。
定义连续的三角形函数
\[
g(x)=\max\left\{0,1-\frac{|x-c|}{\varepsilon}\right\},\qquad x\in[0,1].
\]
它满足$g(c)=1$，所以不恒为零；在$[c-\varepsilon,c+\varepsilon]$外为零，而在该闭区间内$f$为零。因此对所有$x\in[0,1]$均有$f(x)g(x)=0$，即$fg=0$。由于$f\ne0$，$f$为零因子。
\end{proof}
\paragraph{复习提示}
“每个$f\in I$各有零点”与“所有$f\in I$有同一个公共零点”的量词次序不同。第二问不能只要求$f$有一个零点；必须有一段区间。零函数满足零点集条件，但按约定不称零因子，所以题目须限定$f\ne0$。

\subsection{习题44：有限域上的矩阵计数}
\paragraph{题目}
设$p$为素数，$F=\F_p$为$p$元域，$n\geq1$。证明
\[
|M_n(F)|=p^{n^2},\qquad
|\GL_n(F)|=\prod_{k=0}^{n-1}(p^n-p^k).
\]
这里$M_n(F)$是全体$n$阶矩阵的集合，$\GL_n(F)$是其中的可逆矩阵乘法群。
\paragraph{考点与思路}
可逆矩阵的列向量恰为$F^n$的一组有序基。逐列选取时，排除前面列向量张成的子空间。
\textbf{原文依据：}现有笔记本题解答中按矩阵位置计数、再按线性无关列向量逐列计数的证明。
\paragraph{所用结论}
\toolref{g-basis-count}。
\begin{proof}[详细证明]
一个$n$阶矩阵共有$n^2$个位置，每一位置均可独立选取$F$的$p$个元素之一。因此乘法计数原理给出$|M_n(F)|=p^{n^2}$。

现在计数可逆矩阵。线性代数的可逆性判别说明，矩阵可逆当且仅当其$n$个列向量线性无关。第一列不能是零向量，所以有$p^n-1$种选择。

假设已选出线性无关的前$k$列$v_1,\ldots,v_k$，其中$0\leq k<n$。其张成空间中的向量均可写成
\[
\lambda_1v_1+\cdots+\lambda_kv_k,\qquad \lambda_i\in F.
\]
因为这$k$列线性无关，不同系数组给出不同向量：若两组给出相同向量，相减后得到线性关系，所有系数差必须为零。因此该张成空间恰有$p^k$个向量。
第$k+1$列必须位于该空间之外，才能保持列向量线性无关，所以有$p^n-p^k$种选择。

依次选满$n$列，每个有序选择恰对应一个矩阵，没有额外的排列除数。故
\[
|\GL_n(F)|=(p^n-1)(p^n-p)\cdots(p^n-p^{n-1}).
\]
这正是所求乘积。
\end{proof}
\paragraph{复习提示}
当$n=1$时结果为$p-1$，对应域中非零元素的个数。关键区别是：每一列都非零不足以保证可逆，还必须排除前面列向量的整个张成空间。
